\section{Background and Related Work}
\label{sec:background}

\subsection{Smart Grid Co-simulation and Security Datasets}

Smart grid security research relies on cyber-physical testbeds and co-simulation frameworks to study interactions between power system dynamics and control mechanisms \cite{cintuglu2017testbeds,yohanandhan2022review}. HELICS has emerged as a widely used open-source framework for multi-domain co-simulation, supporting time-coordinated interaction between transmission solvers, distribution simulators, and communication models \cite{hardy2024helics}. Security-focused platforms such as GridAttackSim \cite{le2020gridattacksim} combine power system and network simulation but rely on predefined attack scenarios. These platforms represent \emph{Generation~1} evaluation tools: high-fidelity environments with adversary logic fixed in advance.

Similarly, widely used ICS security datasets including SWaT, WADI, EPIC, and HAI \cite{swatdataset,wadidataset,epicdataset,haidataset} capture static recordings of scripted attack campaigns. Survey work highlights that evaluation results from such traces may not generalize to adaptive adversaries \cite{sahani2023smartgridids}. Moreover, most testbeds do not cover EV or DER-rich feeders or load-altering attacks that exploit controllable demand. This motivates evaluation frameworks where dynamic attackers operate \emph{within} co-simulations rather than being pre-scripted around them.

\subsection{Automated Attack Generation and LLM-based Red Teaming}

The security literature has developed models for false data injection attacks (FDIAs) \cite{aoufi2020fdiasurvey} and load-altering attacks \cite{laaSurvey2024}, but these typically assume offline attack planning rather than online adaptation. Automated approaches such as attack graph generation \cite{a2g2v}, control-theoretic red teaming \cite{tackett2024redteam}, and deep reinforcement learning for targeted attacks \cite{targetedAttackSynthesis} represent \emph{Generation~2--3} adversaries that can adapt over time but require scenario-specific state representations.

Recent work applies LLMs to offensive security. PentestGPT \cite{deng2024pentestgpt} and AutoAttacker \cite{autoattacker} demonstrate LLM-based penetration testing in IT systems, while AttackLLM \cite{attackllm} generates attack patterns for ICS testbeds. These \emph{Generation~4} approaches use LLMs as cognitive orchestrators that reason over documentation and call external tools. However, no existing framework (i) couples LLM agents directly into HELICS-style T\&D co-simulation, (ii) exposes physically constrained attack primitives, and (iii) explicitly studies the evaluation validity gap between static and adaptive attackers. LLM-GridEval addresses this gap.